\chapter{Formy przecięcia i uchwyty w wymiarze cztery}
\label{ch:freedman-formy}
\index{forma przecięcia!w wymiarze cztery}

W wymiarze cztery dwuwymiarowe powierzchnie przecinają
się w punktach. Liczby tych przecięć układają się w
symetryczną formę całkowitą. Potrafimy ją obliczać z
homologii, a także z obramowanych diagramów uchwytów.
Pytanie, czy taka forma rozstrzyga typ rozmaitości, jest
już znacznie głębsze. Odpowiedź Freedmana dotyczy
\emph{homeomorfizmów} jednospójnych rozmaitości
czterowymiarowych; nie klasyfikuje ich struktur gładkich.

Przeprowadzimy potrzebną algebrę, w tym pełny rachunek
wyznacznika i sygnatury formy $E_8$. Następnie
obliczymy formy z obramowanych diagramów uchwytów
i wyjaśnimy, dlaczego zwykły gładki trik Whitneya
nie przenosi się automatycznie do wymiaru cztery.
To przygotuje konstrukcje Cassona i topologiczne
twierdzenie Freedmana w następnych rozdziałach.

\section{Forma przecięcia i trzy kategorie}
\label{sec:freedman-form}

Niech $X$ będzie zamkniętą, spójną, zorientowaną
rozmaitością $4$-wymiarową. Forma przecięcia została
wprowadzona w Twierdzeniu~\ref{thm:forma-przeciecia-16}:
\begin{equation}\label{eq:freedman-form}
 Q_X(a,b)=\bigl\langle
  \operatorname{PD}^{-1}(a)\smile
  \operatorname{PD}^{-1}(b),[X]\bigr\rangle,
 \qquad a,b\in H_2(X;\Z)/\operatorname{Tor}.
\end{equation}
Ponieważ stopień obu klas kohomologii jest równy $2$,
przemienność stopniowana daje $Q_X(a,b)=Q_X(b,a)$.
Dualność Poincarégo daje unimodularność: w bazie części
wolnej macierz ma wyznacznik $\pm1$. W kategorii DIFF
dwie poprzeczne zorientowane powierzchnie reprezentujące
$a,b$ mają liczbę przecięcia równą $Q_X(a,b)$, zgodnie
z Twierdzeniem~\ref{thm:przeciecie-kubek-16}.
W kategorii TOP definicją pozostaje powyższe parowanie
homologiczne; opis geometryczny wymaga osobnego twierdzenia
o topologicznej transwersalności.

W przypadku $\pi_1(X)=0$ znika także torsja w $H_2(X;\Z)$.
Istotnie, $H_1(X;\Z)=0$, więc twierdzenie o współczynnikach
uniwersalnych daje
$H^2(X;\Z)\cong\operatorname{Hom}(H_2(X;\Z),\Z)$.
Dualność Poincarégo utożsamia lewą stronę z
$H_2(X;\Z)$. Grupa po prawej jest wolna, zatem
$H_2(X;\Z)$ jest wolna i forma~\eqref{eq:freedman-form}
jest określona na całej grupie. Dla rozmaitości
topologicznych używamy tu topologicznej dualności
Poincarégo, a nie rachunku form różniczkowych.

Ważne są trzy różne znaczenia słowa „rozmaitość”:
\begin{itemize}[leftmargin=2em]
\item w kategorii \emph{TOP} zmiany współrzędnych są
homeomorfizmami; to kategoria twierdzenia Freedmana;
\item w kategorii \emph{PL} mapy przejścia są kawałkami
liniowe po odpowiednim podziale; w wymiarze cztery
struktura PL dopuszcza zgodne wygładzenie;
\item w kategorii \emph{DIFF} zmiany współrzędnych są
gładkie; tylko tutaj stosujemy metryki, koneksje i
gładkie izotopie z poprzednich rozdziałów.
\end{itemize}
Każda struktura gładka daje strukturę PL i topologiczną.
Fakt o wygładzaniu struktur PL w wymiarze cztery jest
wynikiem teorii wygładzania, opisanym przez
Freedmana i Quinna, \emph{Topology of 4-Manifolds},
\S\,8.7; nie wynika z samej definicji PL.
Homeomorfizm dwóch rozmaitości gładkich nie musi być
dyfeomorfizmem. W dalszej części wszystkie twierdzenia
o $Q_X$ dotyczą najpierw kategorii TOP, o ile nie
zaznaczymy inaczej.

\section{Parzystość, unimodularność i sygnatura}
\label{sec:freedman-lattices}

\begin{definicja}[Całkowita forma symetryczna]
\label{def:freedman-integral-form}
Niech $L\cong\Z^r$ będzie wolną grupą abelową skończonej
rangi. Symetryczne odwzorowanie dwuliniowe
$Q:L\times L\to\Z$ nazywamy \emph{formą całkowitą}.
W bazie $e_1,\ldots,e_r$ ma macierz
$A=(Q(e_i,e_j))$. Forma jest \emph{unimodularna},
jeśli $\det A=\pm1$, a \emph{parzysta}, jeśli
$Q(x,x)\in2\Z$ dla każdego $x\in L$;
w przeciwnym razie jest \emph{nieparzysta}.
Jej \emph{sygnatura} $\sigma(Q)$ to liczba dodatnich
wartości własnych $A$ minus liczba ujemnych, liczonych
z krotnościami.
\end{definicja}

Intuicja geometryczna: $Q(x,x)$ mierzy samoprzecięcie
klasy środkowego wymiaru, a parzystość pyta, czy
wszystkie te liczby są parzyste. Unimodularność mówi,
że parowanie rozróżnia każdą całkowitą klasę w jej
dualnej kracie. Sygnatura rejestruje przewagę dodatnich
nad ujemnymi kierunkami po przejściu z $\Z$ do $\R$.

\begin{stwierdzenie}[Rachunek w bazie]
\label{prop:freedman-basis}
Niech $A$ będzie macierzą całkowitej formy symetrycznej.
Wtedy forma jest parzysta wtedy i tylko wtedy, gdy
wszystkie diagonalne wpisy $A_{ii}$ są parzyste.
Przy zmianie bazy przez $P\in\mathrm{GL}_r(\Z)$ macierz
przechodzi w $P^{\mathsf T}AP$; parzystość,
unimodularność i sygnatura nie zmieniają się.
\end{stwierdzenie}
\begin{proof}
Dla $x=\sum_i x_i e_i$ symetria daje
\begin{equation}\label{eq:freedman-parity}
 Q(x,x)=\sum_i A_{ii}x_i^2
             +2\sum_{i<j}A_{ij}x_ix_j.
\end{equation}
Jeśli diagonala jest parzysta, prawa strona jest
parzysta dla każdego $x$. Odwrotnie, podstawienie
$x=e_i$ wykrywa $A_{ii}$. W nowej bazie kolumny
macierzy $P$ są współrzędnymi nowych generatorów,
stąd macierz $P^{\mathsf T}AP$. Jej wyznacznik wynosi
$(\det P)^2\det A=\det A$, bo $\det P=\pm1$.
Parzystość jest własnością wartości $Q(x,x)$,
nie zapisu. Po rozszerzeniu skalarów do $\R$
prawo bezwładności Sylvestera zachowuje liczby
dodatnich i ujemnych kierunków przy kongruencji.
\end{proof}

\begin{przyklad}[Dwie formy o tym samym rzędzie i sygnaturze]
\label{prz:freedman-H-I}
Forma hiperboliczna $H$ i forma diagonalna $I_{1,1}$
mają macierze
\[
 H=\begin{pmatrix}0&1\\1&0\end{pmatrix},
 \qquad
 I_{1,1}=\begin{pmatrix}1&0\\0&-1\end{pmatrix}.
\]
Obie są unimodularne, mają rząd $2$ i sygnaturę $0$:
ich wartości własne są odpowiednio $(1,-1)$ i $(1,-1)$.
Jednak $Q_H((a,b),(a,b))=2ab$ na diagonali kwadratowej,
więc $H$ jest parzysta, natomiast
$Q_{I_{1,1}}((1,0),(1,0))=1$, więc $I_{1,1}$ jest nieparzysta.
Nie mogą być izometryczne nad $\Z$.
Geometria dostarcza przykładów: $Q_{S^2\times S^2}=H$
jak w rozdziale~\ref{ch:przeciecia-dualnosc-klasy},
a dwie linie projektowe w $\C P^2$ przecinają się
raz z dodatnim znakiem, więc $Q_{\C P^2}=[1]$.
Po odwróceniu orientacji otrzymujemy $[-1]$.
Suma spójna $\C P^2\#\overline{\C P^2}$
ma formę $I_{1,1}$: klasy z obu summandów można
reprezentować poza kulami użytymi do sklejenia,
więc ich mieszane przecięcie wynosi zero.
\end{przyklad}

Suma prosta form odpowiada blokowej macierzy diagonalnej.
Z samego rachunku bloków wynikają
$\det(Q\oplus Q')=(\det Q)(\det Q')$ oraz
$\sigma(Q\oplus Q')=\sigma(Q)+\sigma(Q')$.
Nie wyprowadzamy stąd klasyfikacji wszystkich form
całkowitych. Nawet przy ustalonym rzędzie i sygnaturze
typ parzysty i nieparzysty mogą się różnić, a formy
określone w wyższych rangach mają dalsze subtelności.

\section{Macierz formy \texorpdfstring{$E_8$}{E8}}
\label{sec:freedman-E8}

Rozważmy graf będący łańcuchem siedmiu wierzchołków
$1-2-3-4-5-6-7$ z dodatkowym wierzchołkiem $8$
dołączonym do $3$. Umieśćmy $2$ na przekątnej
macierzy, $-1$ przy każdej krawędzi grafu i $0$
w pozostałych miejscach. Otrzymujemy
\begin{equation}\label{eq:freedman-E8}
 E_8=\begin{pmatrix}
 2&-1& 0& 0& 0& 0& 0& 0\\
 -1& 2&-1& 0& 0& 0& 0& 0\\
 0&-1& 2&-1& 0& 0& 0&-1\\
 0& 0&-1& 2&-1& 0& 0& 0\\
 0& 0& 0&-1& 2&-1& 0& 0\\
 0& 0& 0& 0&-1& 2&-1& 0\\
 0& 0& 0& 0& 0&-1& 2& 0\\
 0& 0&-1& 0& 0& 0& 0& 2
 \end{pmatrix}.
\end{equation}
W literaturze spotyka się również macierz z $+1$
przy krawędziach. Ponieważ graf jest drzewem, można
zmieniać znaki wybranych wektorów bazy od jednego
wierzchołka do następnego i otrzymać tę drugą macierz;
obie zapisują tę samą formę całkowitą.

\begin{figure}[htbp]
\centering
\begin{tikzpicture}[font=\small,>=Latex]
 \foreach \i in {1,...,7}{
  \node[circle,draw=manifoldblue!80!black,
        fill=manifoldblue!12,minimum size=8mm,inner sep=0pt]
       (v\i) at ({1.12*(\i-1)},0) {$2$};
  \node[below=2mm of v\i] {$e_{\i}$};
 }
 \node[circle,draw=manifoldblue!80!black,
       fill=manifoldblue!12,minimum size=8mm,inner sep=0pt]
       (v8) at (2.24,1.25) {$2$};
 \node[above=2mm of v8] {$e_8$};
 \foreach \i/\j in {1/2,2/3,3/4,4/5,5/6,6/7,3/8}{
  \draw[manifoldblue!80!black,thick] (v\i)--(v\j);
 }
 \node[anchor=west,align=left] at (7.15,.65)
   {krawędź: $Q(e_i,e_j)=-1$\\brak krawędzi: $Q(e_i,e_j)=0$};
\end{tikzpicture}
\caption{Graf obramowań dla formy $E_8$.
Każdy wierzchołek oznacza uchwyt $2$ o obramowaniu $+2$,
a krawędź pojedyncze sprzężenie z liczbą $-1$.
Graf koduje macierz przecięć; nie jest rzutem całego
ośmioskładnikowego splotu na płaszczyznę.}
\label{fig:freedman-E8-graph}
\end{figure}

\begin{stwierdzenie}[Parzystość, wyznacznik i sygnatura $E_8$]
\label{prop:freedman-E8-properties}
Forma~\eqref{eq:freedman-E8} jest parzysta,
unimodularna i dodatnio określona. Dokładniej,
\[
 \det E_8=1,\qquad \sigma(E_8)=8.
\]
\end{stwierdzenie}
\begin{proof}
Parzystość wynika już z diagonalnych wpisów $2$
i Stwierdzenia~\ref{prop:freedman-basis}. Jawnie,
jeśli $x=(x_1,\ldots,x_8)\in\Z^8$, to
\[
 x^{\mathsf T}E_8x
 =2\sum_{i=1}^8x_i^2
       -2\!\sum_{\{i,j\}\,\text{jest krawędzią}}x_ix_j
 \in2\Z.
\]

Niech $A_m$ oznacza macierz łańcucha $m$ wierzchołków:
$2$ na przekątnej, $-1$ bezpośrednio nad nią i pod nią.
Rozwinięcie wyznacznika według ostatniego wiersza daje
$d_m:=\det A_m=2d_{m-1}-d_{m-2}$, przy
$d_0=1$, $d_1=2$. Indukcja daje $d_m=m+1$.
W szczególności $\det A_7=8$ i wszystkie jego wiodące
minory są dodatnie; kryterium Sylvestera mówi,
że $A_7$ jest dodatnio określona.

Macierz $E_8$ zapisujemy blokowo jako
\[
 E_8=\begin{pmatrix}A_7&-e_3\\-e_3^{\mathsf T}&2\end{pmatrix},
\]
gdzie $e_3$ jest trzecim wektorem standardowym
w $\R^7$. Usunięcie trzeciego wiersza i kolumny
z $A_7$ pozostawia bloki $A_2$ i $A_4$.
Reguła Cramera daje więc
\[
 (A_7^{-1})_{33}
 =\frac{\det A_2\det A_4}{\det A_7}
 =\frac{3\cdot5}{8}=\frac{15}{8}.
\]
Dopełnienie Schura ostatniego wpisu ma wartość
$2-e_3^{\mathsf T}A_7^{-1}e_3=1/8>0$.
Po kongruencji blokowej eliminującej wyrazy
pozadiagonalne otrzymujemy
$A_7\oplus[1/8]$. Stąd $E_8$ jest dodatnio
określona, ma osiem dodatnich kierunków, a
\[
 \det E_8=\det A_7\cdot\frac18=8\cdot\frac18=1.
\]
To obliczenie dowodzi wszystkich trzech własności
bez odwołania do klasyfikacji kratek.
\end{proof}

Forma $-E_8$ jest również parzysta i unimodularna,
lecz ma sygnaturę $-8$. Formy $E_8\oplus H$,
$E_8\oplus E_8$ i $H\oplus H$ dają następne
bezpośrednio sprawdzalne przykłady. Nie twierdzimy,
że każda forma parzysta unimodularna jest taką sumą.
Ogólne twierdzenie algebraiczne o podzielności
sygnatury \emph{dowolnej} parzystej formy unimodularnej
przez $8$ oraz klasyfikacja form nieokreślonych
znajdują się w \href{https://archive.mpim-bonn.mpg.de/4789/2/FreedmanQuinn-TopologyOf4Manifolds-Reformatted2013.pdf}
{Freedman--Quinn, \emph{Topology of 4-Manifolds},
\S\,10.2.1} (z odesłaniem do Milnora--Husemollera).
Przyjmujemy je jako wyniki zewnętrzne tylko wtedy,
gdy formułujemy pełne twierdzenie Freedmana;
nie używamy ich do powyższego rachunku $E_8$.

\section{Obramowane uchwyty i granica triku Whitneya}
\label{sec:freedman-handles}

Rozdział~\ref{ch:rachunek-uchwytow} wprowadził uchwyt
$k$ jako $D^k\times D^{4-k}$ przyczepiony wzdłuż
$S^{k-1}\times D^{4-k}$. Dla uchwytu $2$ w wymiarze
cztery krzywą przyczepienia jest knot
$K\subset\partial D^4=S^3$. Obramowanie jego wiązki
normalnej w $S^3$ podaje, ile razy pas $K\times D^2$
skręca względem obramowania zerowego od powierzchni
Seiferta. Tę różnicę mierzy liczba całkowita $f$.

\begin{definicja}[Obramowany diagram uchwytów $2$]
\label{def:freedman-framed-link}
\emph{Obramowany splot} $L=K_1\cup\cdots\cup K_r$
w $S^3$ to rozłączne zorientowane knoty z liczbami
całkowitymi $f_1,\ldots,f_r$. Oznacza zwartą gładką
rozmaitość z brzegiem
\[
 W_L=D^4\cup_{K_1,f_1}(D^2\times D^2)
       \cup\cdots\cup_{K_r,f_r}(D^2\times D^2).
\]
Jej \emph{macierz sprzężeń} $B_L$ ma
$(B_L)_{ii}=f_i$ oraz
$(B_L)_{ij}=\operatorname{lk}(K_i,K_j)$ dla $i\ne j$.
\end{definicja}

Liczba obramowania ma znaczenie geometryczne:
odpowiada samoprzecięciu powierzchni zbudowanej z rdzenia
uchwytu i powierzchni rozpinającej knot. Sprzężenie
dwóch komponentów odpowiada przecięciu dwóch takich
powierzchni. Ten opis pozwala przejść od rysunku w $S^3$
do macierzy formy w rozmaitości $4$-wymiarowej.

\begin{stwierdzenie}[Macierz uchwytów $2$]
\label{prop:freedman-link-matrix}
Dla $W_L$ z Definicji~\ref{def:freedman-framed-link}
grupa $H_2(W_L;\Z)$ jest wolna rangi $r$, a jej
forma przecięcia w bazie wyznaczonej przez rdzenie
uchwytów domknięte powierzchniami Seiferta ma macierz $B_L$.
Jeżeli $\det B_L=\pm1$, to $\partial W_L$ jest
całkowitą sferą homologiczną wymiaru trzy.
\end{stwierdzenie}
\begin{proof}
Względny kompleks uchwytowy pary $(W_L,D^4)$ ma
$r$ generatorów w stopniu $2$ i nie ma generatorów
w stopniach $1$ ani $3$. Stąd
$H_2(W_L;\Z)\cong\Z^r$ i $H_1(W_L;\Z)=0$.
Do rdzenia każdego uchwytu doklejamy powierzchnię
Seiferta komponentu $K_i$, przesuniętą nieco do
wnętrza $D^4$. Otrzymane zamknięte powierzchnie
reprezentują bazę $H_2(W_L;\Z)$.
Po poprzecznym odsunięciu jednej z nich liczba
przecięcia z drugą jest liczbą sprzężenia
$\operatorname{lk}(K_i,K_j)$. Przy $i=j$ odsunięcie
wykorzystuje obramowanie, a liczba przecięcia jest
$f_i$ względem obramowania Seiferta. To daje $B_L$.

Niech $Y=\partial W_L$. Dualność Lefschetza i
twierdzenie o współczynnikach uniwersalnych
utożsamiają $H_2(W_L,Y;\Z)$ z
$\operatorname{Hom}(H_2(W_L;\Z),\Z)$, ponieważ
$H_1(W_L)=0$. W ciągu dokładnym pary odwzorowanie
$H_2(W_L)\to H_2(W_L,Y)$ jest reprezentowane przez
$B_L$. Gdy $\det B_L=\pm1$, jest izomorfizmem.
Wtedy dokładność daje $H_1(Y)=0$.
Ponadto $H_3(W_L,Y)\cong H^1(W_L)=0$, więc również
$H_2(Y)=0$. Brzeg $Y$ jest spójną orientowalną
zamkniętą rozmaitością $3$-wymiarową, zatem ma
homologie $S^3$.
\end{proof}

Graf z Rysunku~\ref{fig:freedman-E8-graph} można
zrealizować jako obramowany splot: każdy wierzchołek
zastępujemy niezawęźlonym okręgiem z obramowaniem
$+2$, a dla każdej krawędzi wprowadzamy pojedyncze
sprzężenie Hopfa o znaku $-1$. Wtedy macierz
$B_L$ to~\eqref{eq:freedman-E8}. Powyższe
stwierdzenie dowodzi, że brzeg tak zbudowanego
gładkiego $W_L$ jest sferą \emph{homologiczną},
nie że jest standardowym $S^3$. Domknięcie tej
rozmaitości przez kontraktywne \emph{topologiczne}
wypełnienie brzegu jest osobnym głębokim krokiem
Freedmana; omawiają go
\href{https://archive.mpim-bonn.mpg.de/4789/2/FreedmanQuinn-TopologyOf4Manifolds-Reformatted2013.pdf}
{Freedman--Quinn, \S\,9.7 i \S\,10.5.1}.

W dowodzie wysokowymiarowego twierdzenia o
$h$-kobordyzmie dwa punkty przecięcia o przeciwnych
znakach usuwamy dyskiem Whitneya
(Sekcja~\ref{sec:whitney-hc-19}). W wymiarze cztery
zarówno powierzchnie, jak i dysk Whitneya mają
wymiar $2$. Transwersalność przewiduje
\begin{equation}\label{eq:freedman-whitney-count}
 \dim(D^2\cap\Sigma^2)_{\mathrm{oczek.}}
       =2+2-4=0.
\end{equation}
Wnętrze dysku może więc przeciąć powierzchnię
w kolejnych izolowanych punktach. Taki sam rachunek
dotyczy samoprzecięć dysku. W wymiarze otaczającym
co najmniej $5$ oczekiwany wymiar jest ujemny i
transwersalne odsunięcie usuwa te przecięcia.
Sama równość sumy znaków punktów z zerem nie dostarcza
w wymiarze cztery zanurzonego, obramowanego dysku
o czystym wnętrzu.

\begin{figure}[htbp]
\centering
\begin{tikzpicture}[font=\small,>=Latex]
 \node[draw=manifoldblue!70!black,rounded corners,
       minimum width=5.7cm,minimum height=3.3cm] at (-3.2,0) {};
 \node[font=\bfseries] at (-3.2,1.05)
   {dysk Whitneya};
 \draw[manifoldblue!80!black,very thick]
   (-5.65,-.55)--(-.75,-.55);
 \draw[manifoldred!85!black,very thick]
   (-5.45,-1.20) .. controls (-4.85,.42) and (-1.55,.42)
       .. (-.95,-1.20);
 \fill[manifoldgold!24]
   (-4.95,-.55) .. controls (-4.25,.08) and (-2.15,.08)
       .. (-1.45,-.55)--cycle;
 \draw[manifoldgold!85!black,thick]
   (-4.95,-.55) .. controls (-4.25,.08) and (-2.15,.08)
       .. (-1.45,-.55);
 \fill[manifoldblue!80!black] (-4.95,-.55) circle (2pt)
       node[below] {$p_+$};
 \fill[manifoldblue!80!black] (-1.45,-.55) circle (2pt)
       node[below] {$p_-$};
 \fill[manifoldred!85!black] (-3.15,-.25) circle (2pt)
       node[above] {$q$};
 \node[manifoldgold!70!black] at (-3.25,-.83) {$D$};
 \node[anchor=west,align=left] at (.15,.65)
   {$D^2\cap\Sigma^2$ w $4$ wymiarach:\\
      $2+2-4=0$ --- punkty $q$ mogą pozostać};
 \node[anchor=west,align=left] at (.15,-.55)
   {$D^2\cap\Sigma^2$ w $n\geq5$:\\
      $2+2-n<0$ --- transwersalnie pusto};
\end{tikzpicture}
\caption{Rysunek jest jedynie rzutem schematycznym:
łuki dwóch powierzchni łączą punkty $p_+$ i $p_-$,
a złoty dysk $D$ może mieć dodatkowe przecięcie $q$
we wnętrzu. Rachunek wymiarów wyjaśnia, dlaczego
transwersalne odsunięcie usuwa takie przecięcia
w wyższych wymiarach, ale nie rozstrzyga tego w wymiarze $4$.}
\label{fig:freedman-whitney}
\end{figure}

To jest miejsce, w którym potrzebne będą uchwyty
Cassona i topologiczne twierdzenie o zanurzaniu dysku.
Nie wyprowadzamy ich z rachunku wymiarów; ten rachunek
wyjaśnia tylko, dlaczego znany gładki dowód nie działa.

\par\smallskip
{\footnotesize\noindent\emph{Dalsza lektura:}
\href{https://archive.mpim-bonn.mpg.de/4789/2/FreedmanQuinn-TopologyOf4Manifolds-Reformatted2013.pdf}
{M.\ H.\ Freedman i F.\ Quinn, \emph{Topology of
4-Manifolds} (1990)}, zwłaszcza rozdziały 9--10:
forma przecięcia, uchwyty i geometryczne trudności
w wymiarze cztery.\par}
